\section{Conclusions}\label{sec:conclusions}

\subsection{Summary of Findings}
As long as cellular connections are maintained, \textbf{a high re-identification risk} persists due to behavioral predictability \cite{duan2025sure, duan2025moco}. The analysis presented in this report leads to the following key conclusions:

\begin{itemize}
    \item \textbf{Behavioral metadata can function as a highly distinctive identifier.} The SURE framework demonstrates that BS association patterns and traffic usage create unique behavioral fingerprints, achieving re-identification AUC scores of 0.941 on datasets of 10,000 users, even with k-anonymity applied \cite{duan2025sure}.
    \item \textbf{Social graphs are exposed.} MoCo extends the threat from individual identification to social network reconstruction, detecting co-traveling contacts with an F1-score of 87.13\% using only cellular signaling traces \cite{duan2025moco}.
    \item \textbf{Software mitigations are insufficient.} Traditional anonymization techniques and model-hardening approaches fail to address the root cause: centralized metadata collection by cellular operators. While defenses such as ASTRA \cite{chhipa2025astra} and RoCL \cite{kim2020rocl} improve model robustness, they operate at the learning level and leave the underlying data pipeline untouched.
    \item \textbf{Physical-layer separation offers a compelling alternative.} The Meshtastic parallel radio network eliminates the cellular operator metadata source exploited by both SURE and MoCo by operating on a completely independent infrastructure with no centralized logging \cite{meshtastic_overview}.
\end{itemize}

Parallel radio networks like Meshtastic offer a pathway for operating outside the centralized cellular data ecosystem \cite{meshtastic_overview}.

\subsection{Future Extensions}
Several directions for future work emerge from this analysis:

\begin{itemize}
    \item \textbf{Multi-hop scalability.} Investigating the performance of Meshtastic mesh networks at scale, particularly the impact of increasing hop counts on latency, reliability, and metadata leakage through traffic analysis.
    \item \textbf{Hybrid encrypted mesh systems.} Combining Meshtastic's LoRa backbone with short-range protocols (e.g., BLE mesh) to increase bandwidth for local communications while maintaining the metadata-free properties of the long-range network.
    \item \textbf{RF fingerprinting resistance.} Developing countermeasures against emerging RF fingerprinting techniques that could potentially identify individual LoRa transmitters based on hardware-specific signal characteristics.
    \item \textbf{Experimental validation.} Conducting the proposed evaluation framework with real user groups to empirically quantify the re-identification risk reduction achieved by Meshtastic compared to cellular baselines.
\end{itemize}
\newpage